\documentclass[../../main.tex]{subfiles}
\begin{document}
\section{斜对称张量的代数}\label{section:ch14-alternating-tensors}

设 $V$ 是一个有限维（实）向量空间。回忆一下，$V$ 上的一个协变 $k$-张量，只要交换任意两个自变量其值就变号，就称为\emph{斜对称的}；等价地说，只要自变量的任意排列都使其值乘以该排列的符号（见习题\ref{exercise:ch12-17}）。斜对称协变 $k$-张量也称为\emph{外形式}、\emph{多重余向量}或 $k$-余向量。$V$ 上全体 $k$-余向量所成的向量空间记作 $\Lambda^k(V^*)$。正如我们在第~12 章中注意到的，所有 $0$-张量与 $1$-张量都是斜对称的。

下一个引理给出斜对称张量另外两个有用的刻画。

\begin{lemma}\label{lemma:ch14-alternating-equiv}
设 $\alpha$ 是有限维向量空间 $V$ 上的一个协变 $k$-张量。下列各条等价：
\begin{enumerate}[(a)]
\item $\alpha$ 是斜对称的。
\item 只要 $k$ 元组 $(v_1,\dots,v_k)$ 线性相关，就有 $\alpha(v_1,\dots,v_k)=0$。
\item $\alpha$ 在其任意两个自变量相等时取零值：
\begin{align*}
\alpha(v_1,\dots,w,\dots,w,\dots,v_k)=0.
\end{align*}
\end{enumerate}
\end{lemma}

\begin{proof}
蕴含关系 (a)$\Rightarrow$(c) 与 (b)$\Rightarrow$(c) 都是显然的。我们通过证明 (c) 同时蕴含 (a) 与 (b) 来完成证明。

假设 $\alpha$ 满足 (c)。对任意向量 $v_1,\dots,v_k$，该假设蕴含
\begin{align*}
0&=\alpha(v_1,\dots,v_i+v_j,\dots,v_i+v_j,\dots,v_k)\\
&=\alpha(v_1,\dots,v_i,\dots,v_i,\dots,v_k)+\alpha(v_1,\dots,v_i,\dots,v_j,\dots,v_k)\\
&\quad+\alpha(v_1,\dots,v_j,\dots,v_i,\dots,v_k)+\alpha(v_1,\dots,v_j,\dots,v_j,\dots,v_k)\\
&=\alpha(v_1,\dots,v_i,\dots,v_j,\dots,v_k)+\alpha(v_1,\dots,v_j,\dots,v_i,\dots,v_k).
\end{align*}
于是 $\alpha$ 是斜对称的。另一方面，若 $(v_1,\dots,v_k)$ 是一个线性相关的 $k$ 元组，则其中某个 $v_i$ 可写成其余向量的线性组合。为简单起见，设 $v_k=\sum_{j=1}^{k-1}a_jv_j$。由 $\alpha$ 的多重线性性，
\begin{align*}
\alpha(v_1,\dots,v_k)=\sum_{j=1}^{k-1}a_j\,\alpha(v_1,\dots,v_{k-1},v_j).
\end{align*}
在每一项中 $\alpha$ 都有两个相同的自变量，故每一项都为零。
\end{proof}

在第~12 章中，我们定义了一个称为对称化的投影 $\operatorname{Sym}:T^k(V^*)\to S^k(V^*)$。我们定义一个类似的投影 $\operatorname{Alt}:T^k(V^*)\to\Lambda^k(V^*)$，称为\emph{斜对称化}，如下：
\begin{align*}
\operatorname{Alt}\alpha=\frac1{k!}\sum_{\sigma\in S_k}(\operatorname{sgn}\sigma)\,{}^\sigma\alpha,
\end{align*}
其中 $S_k$ 是 $k$ 个元素上的对称群。更明确地说，这意味着
\begin{align*}
(\operatorname{Alt}\alpha)(v_1,\dots,v_k)
=\frac1{k!}\sum_{\sigma\in S_k}(\operatorname{sgn}\sigma)\,
\alpha\bigl(v_{\sigma(1)},\dots,v_{\sigma(k)}\bigr).
\end{align*}

\begin{example}\label{example:ch14-alt}
若 $\alpha$ 是任意 $1$-张量，则 $\operatorname{Alt}\alpha=\alpha$。若 $\beta$ 是一个 $2$-张量，则
\begin{align*}
(\operatorname{Alt}\beta)(v,w)=\frac12\bigl(\beta(v,w)-\beta(w,v)\bigr).
\end{align*}
对一个 $3$-张量 $\gamma$，
\begin{align*}
(\operatorname{Alt}\gamma)(v,w,x)=\frac16\bigl(&\gamma(v,w,x)+\gamma(w,x,v)+\gamma(x,v,w)\\
&-\gamma(w,v,x)-\gamma(v,x,w)-\gamma(x,w,v)\bigr).
\end{align*}
\end{example}

下一个命题是命题\ref{proposition:ch12-symmetrization}在斜对称张量上的类比。

\begin{proposition}[斜对称化的性质]\label{proposition:ch14-alt-props}
设 $\alpha$ 是有限维向量空间上的一个协变张量。
\begin{enumerate}[(a)]
\item $\operatorname{Alt}\alpha$ 是斜对称的。
\item $\operatorname{Alt}\alpha=\alpha$ 当且仅当 $\alpha$ 是斜对称的。
\end{enumerate}
\end{proposition}

\begin{exercise}\label{exercise:ch14-04}
证明命题\ref{proposition:ch14-alt-props}。
\end{exercise}

\subsection{初等斜对称张量}

为便于计算斜对称张量，下面这套记号极为有用。给定一个正整数 $k$，正整数的有序 $k$ 元组 $I=(i_1,\dots,i_k)$ 称为一个长度为 $k$ 的\emph{多重指标}。若 $I$ 是这样一个多重指标，$\sigma\in S_k$ 是 $\{1,\dots,k\}$ 的一个排列，我们记
\begin{align*}
I\sigma=\bigl(i_{\sigma(1)},\dots,i_{\sigma(k)}\bigr).
\end{align*}

设 $V$ 是一个 $n$ 维向量空间，并设 $(\varepsilon^1,\dots,\varepsilon^n)$ 是 $V^*$ 的任意一组基。现在我们在 $V$ 上定义一族推广 $\mathbb{R}^n$ 上行列式函数的 $k$-余向量。对每个长度为 $k$、满足 $1\leqslant i_1,\dots,i_k\leqslant n$ 的多重指标 $I=(i_1,\dots,i_k)$，定义一个协变 $k$-张量 $\varepsilon^I=\varepsilon^{i_1\dots i_k}$：
\begin{align}
\varepsilon^I(v_1,\dots,v_k)
=\det\begin{pmatrix}
\varepsilon^{i_1}(v_1)&\cdots&\varepsilon^{i_1}(v_k)\\
\vdots&&\vdots\\
\varepsilon^{i_k}(v_1)&\cdots&\varepsilon^{i_k}(v_k)
\end{pmatrix}
=\det\begin{pmatrix}
v^{i_1}_1&\cdots&v^{i_1}_k\\
\vdots&&\vdots\\
v^{i_k}_1&\cdots&v^{i_k}_k
\end{pmatrix}.\label{align:ch14-elem}
\end{align}
换言之，若 $v$ 表示以 $v_1,\dots,v_k$ 在与 $(\varepsilon^i)$ 对偶的基 $(E_i)$ 下的分量为列的 $n\times k$ 矩阵，则 $\varepsilon^I(v_1,\dots,v_k)$ 是由 $v$ 的第 $i_1,\dots,i_k$ 行所成的 $k\times k$ 子矩阵的行列式。由于交换任意两列都使行列式变号，$\varepsilon^I$ 显然是一个斜对称 $k$-张量。我们称 $\varepsilon^I$ 为一个\emph{初等斜对称张量}或\emph{初等 $k$-余向量}。

\begin{example}\label{example:ch14-elementary}
在 $\bigl(\mathbb{R}^3\bigr)^*$ 的标准对偶基 $(e^1,e^2,e^3)$ 下，我们有
\begin{align*}
e^{13}(v,w)=v^1w^3-w^1v^3;\qquad
e^{123}(v,w,x)=\det(v,w,x).
\end{align*}
\end{example}

为简化初等 $k$-余向量的计算，把 Kronecker 符号记号作如下推广是有用的。若 $I,J$ 是长度同为 $k$ 的多重指标，我们定义
\begin{align*}
\delta^I_J=\det\begin{pmatrix}
\delta^{i_1}_{j_1}&\cdots&\delta^{i_1}_{j_k}\\
\vdots&&\vdots\\
\delta^{i_k}_{j_1}&\cdots&\delta^{i_k}_{j_k}
\end{pmatrix}.
\end{align*}

\begin{exercise}\label{exercise:ch14-06}
证明
\begin{align*}
\delta^I_J=
\begin{cases}
\operatorname{sgn}\sigma,&\text{若 }I,J\text{ 都无重复指标且 }J=I\sigma\text{ 对某个 }\sigma\in S_k;\\
0,&\text{若 }I\text{ 或 }J\text{ 有重复指标，或 }J\text{ 不是 }I\text{ 的排列}.
\end{cases}
\end{align*}
\end{exercise}

\begin{lemma}[初等 $k$-余向量的性质]\label{lemma:ch14-elementary-props}
设 $(E_i)$ 是 $V$ 的一组基，$(\varepsilon^i)$ 是 $V^*$ 的对偶基，$\varepsilon^I$ 如上定义。
\begin{enumerate}[(a)]
\item 若 $I$ 有重复指标，则 $\varepsilon^I=0$。
\item 若对某个 $\sigma\in S_k$ 有 $J=I\sigma$，则 $\varepsilon^I=(\operatorname{sgn}\sigma)\varepsilon^J$。
\item 在一组基向量上计算 $\varepsilon^I$ 的结果为
\begin{align*}
\varepsilon^I(E_{j_1},\dots,E_{j_k})=\delta^I_J.
\end{align*}
\end{enumerate}
\end{lemma}

\begin{proof}
若 $I$ 有重复指标，则对任意向量 $v_1,\dots,v_k$，\eqref{align:ch14-elem}中的行列式有两行相同，故等于零，这就证明了 (a)。另一方面，若 $J$ 是由 $I$ 交换两个指标所得，则相应的行列式符号相反；这就蕴含 (b)。最后，(c) 直接来自 $\varepsilon^I$ 的定义。
\end{proof}

初等 $k$-余向量的重要性在于，它们为 $\Lambda^k(V^*)$ 提供了一组方便的基。当然，这些 $\varepsilon^I$ 并非全部线性无关，因为其中一些为零，而对应于同一多重指标不同排列的那些彼此成数量倍。但正如下一个命题所表明的，通过只取多重指标的一个适当子集，我们便可得到一组基。

一个多重指标 $I=(i_1,\dots,i_k)$ 若满足 $i_1<\cdots<i_k$，就称为\emph{递增的}。使用带撇的求和号来表示只对递增多重指标求和是方便的，例如
\begin{align*}
\sideset{}{'}\sum_I \alpha_I\varepsilon^I
=\sum_{\{I:i_1<\cdots<i_k\}}\alpha_I\varepsilon^I.
\end{align*}

\begin{proposition}[$\Lambda^k(V^*)$ 的一组基]\label{proposition:ch14-alternating-basis}
设 $V$ 是一个 $n$ 维向量空间。若 $(\varepsilon^i)$ 是 $V^*$ 的任意一组基，则对每个正整数 $k\leqslant n$，$k$-余向量所成之集
\begin{align*}
\mathcal E=\{\varepsilon^I:I\text{ 是长度为 }k\text{ 的递增多重指标}\}
\end{align*}
是 $\Lambda^k(V^*)$ 的一组基。因此
\begin{align*}
\dim\Lambda^k(V^*)=\binom{n}{k}=\frac{n!}{k!(n-k)!}.
\end{align*}
若 $k>n$，则 $\dim\Lambda^k(V^*)=0$。
\end{proposition}

\begin{proof}
当 $k>n$ 时 $\Lambda^k(V^*)$ 是平凡向量空间，这直接由引理\ref{lemma:ch14-alternating-equiv}(b) 得到，因为在那种情形下任意 $k$ 个向量都线性相关。对于 $k\leqslant n$ 的情形，我们需证明 $\mathcal E$ 张成 $\Lambda^k(V^*)$ 且线性无关。设 $(E_i)$ 是 $V$ 的与 $(\varepsilon^i)$ 对偶的基。

为证 $\mathcal E$ 张成 $\Lambda^k(V^*)$，任取 $\alpha\in\Lambda^k(V^*)$。对每个多重指标 $I=(i_1,\dots,i_k)$（不必递增），定义实数 $\alpha_I=\alpha(E_{i_1},\dots,E_{i_k})$。由 $\alpha$ 斜对称可知，若 $I$ 含重复指标则 $\alpha_I=0$；而若 $J=I\sigma$ 对 $\sigma\in S_k$，则 $\alpha_J=(\operatorname{sgn}\sigma)\alpha_I$。于是对任意多重指标 $J$，引理\ref{lemma:ch14-elementary-props}给出
\begin{align*}
\sideset{}{'}\sum_I\alpha_I\varepsilon^I(E_{j_1},\dots,E_{j_k})
=\sideset{}{'}\sum_I\alpha_I\delta^I_J=\alpha_J=\alpha(E_{j_1},\dots,E_{j_k}).
\end{align*}
故 $\sum'_I\alpha_I\varepsilon^I=\alpha$，从而 $\mathcal E$ 张成 $\Lambda^k(V^*)$。

为证 $\mathcal E$ 线性无关，假设对某些系数 $\alpha_I$ 有 $\sum'_I\alpha_I\varepsilon^I=0$。令 $J$ 为任意递增多重指标。把恒等式两边同时作用于向量 $(E_{j_1},\dots,E_{j_k})$，并利用引理\ref{lemma:ch14-elementary-props}，得
\begin{align*}
0=\sideset{}{'}\sum_I\alpha_I\varepsilon^I(E_{j_1},\dots,E_{j_k})=\alpha_J.
\end{align*}
故每个系数 $\alpha_J$ 都为零。
\end{proof}

特别地，对一个 $n$ 维向量空间 $V$，这个命题蕴含 $\Lambda^n(V^*)$ 是一维的，且由 $\varepsilon^{1\dots n}$ 张成。由定义，这个初等 $n$-余向量作用于向量 $(v_1,\dots,v_n)$ 时取其分量矩阵 $v=(v^i_j)$ 的行列式。例如，在 $\mathbb{R}^n$ 上带标准基时，$e^{1\dots n}$ 恰是行列式函数。

由此得到下面这个有用的结论，它刻画了 $n$ 维空间上一个 $n$-余向量在线性映射下的行为。

\begin{proposition}\label{proposition:ch14-n-covector-det}
设 $V$ 是一个 $n$ 维向量空间，$\omega\in\Lambda^n(V^*)$。若 $T:V\to V$ 是任意线性映射，$v_1,\dots,v_n$ 是 $V$ 中任意向量，则
\begin{align}
\omega(Tv_1,\dots,Tv_n)=(\det T)\,\omega(v_1,\dots,v_n).\label{align:ch14-det-pullback}
\end{align}
\end{proposition}

\begin{proof}
设 $(E_i)$ 是 $V$ 的任意一组基，$(\varepsilon^i)$ 是对偶基。设 $(T^j_i)$ 表示 $T$ 关于这组基的矩阵，令 $T_i=TE_i=T^j_iE_j$。由命题\ref{proposition:ch14-alternating-basis}，可写 $\omega=c\,\varepsilon^{1\dots n}$，其中 $c$ 为某个实数。由于\eqref{align:ch14-det-pullback}两边都是 $v_1,\dots,v_n$ 的多重线性函数，只需在 $v_i$ 都是基向量时验证该恒等式。又由于两边都斜对称，只需检查情形 $(v_1,\dots,v_n)=(E_1,\dots,E_n)$。此时\eqref{align:ch14-det-pullback}的右端为
\begin{align*}
(\det T)\,c\,\varepsilon^{1\dots n}(E_1,\dots,E_n)=c\det T.
\end{align*}
另一方面，左端化为
\begin{align*}
\varepsilon^{1\dots n}(TE_1,\dots,TE_n)
=c\,\varepsilon^{1\dots n}(T^j_1E_j,\dots,T^j_nE_j)
=c\det\bigl(T^j_i\bigr),
\end{align*}
它等于右端。
\end{proof}

\subsection{外积}

在第~12 章中，我们定义了对称乘积，它把一对对称张量 $\alpha,\beta$ 变成另一个对称张量 $\alpha\beta=\operatorname{Sym}(\alpha\otimes\beta)$，其秩为原来两个张量秩之和。

本节中我们为斜对称张量定义一个类似的乘积运算。一种定义方式是模仿对称情形，把斜对称张量 $\omega$ 与 $\eta$ 的乘积定义为 $\operatorname{Alt}(\omega\otimes\eta)$。然而我们将使用一个乍看更复杂、但实际上更便于计算的定义。

在本节中我们继续假设 $V$ 是一个有限维实向量空间。给定 $\omega\in\Lambda^k(V^*)$、$\eta\in\Lambda^l(V^*)$，我们定义它们的\emph{外积}或\emph{楔积}为如下 $(k+l)$-余向量：
\begin{align}
\omega\wedge\eta=\frac{(k+l)!}{k!\,l!}\operatorname{Alt}(\omega\otimes\eta).\label{align:ch14-wedge-def}
\end{align}
这个看似神秘的系数，其动机在于下面引理陈述的简洁性。

\begin{lemma}\label{lemma:ch14-wedge-concat}
设 $V$ 是一个 $n$ 维向量空间，$(\varepsilon^1,\dots,\varepsilon^n)$ 是 $V^*$ 的一组基。对任意多重指标 $I=(i_1,\dots,i_k)$ 与 $J=(j_1,\dots,j_l)$，
\begin{align}
\varepsilon^I\wedge\varepsilon^J=\varepsilon^{IJ},\label{align:ch14-wedge-concat}
\end{align}
其中 $IJ=(i_1,\dots,i_k,j_1,\dots,j_l)$ 是把 $I$ 与 $J$ 连接起来所得的多重指标。
\end{lemma}

\begin{proof}
由多重线性性，只需证明对任意基向量序列 $(E_{p_1},\dots,E_{p_{k+l}})$，都有
\begin{align}
\varepsilon^I\wedge\varepsilon^J(E_{p_1},\dots,E_{p_{k+l}})
=\varepsilon^{IJ}(E_{p_1},\dots,E_{p_{k+l}}).\label{align:ch14-wedge-concat-eval}
\end{align}
我们分几种情形考虑。

\emph{情形 1}：$P=(p_1,\dots,p_{k+l})$ 有重复指标。此时由引理\ref{lemma:ch14-alternating-equiv}(c)，\eqref{align:ch14-wedge-concat-eval}两边都为零。

\emph{情形 2}：$P$ 含一个既不出现在 $I$ 中也不出现在 $J$ 中的指标。此时由引理\ref{lemma:ch14-elementary-props}(c)，右端为零。类似地，左端展开中的每一项都涉及 $\varepsilon^I$ 或 $\varepsilon^J$ 在一个并非 $I$ 或 $J$ 排列的基向量序列上取值，故左端也为零。

\emph{情形 3}：$P=IJ$ 且 $P$ 无重复指标。此时由引理\ref{lemma:ch14-elementary-props}(c)，\eqref{align:ch14-wedge-concat-eval}的右端等于 $1$，故我们只需证明左端也等于 $1$。由定义，
\begin{align*}
\varepsilon^I\wedge\varepsilon^J(E_{p_1},\dots,E_{p_{k+l}})
&=\frac{(k+l)!}{k!\,l!}\operatorname{Alt}(\varepsilon^I\otimes\varepsilon^J)(E_{p_1},\dots,E_{p_{k+l}})\\
&=\frac1{k!\,l!}\sum_{\sigma\in S_{k+l}}(\operatorname{sgn}\sigma)\,
\varepsilon^I(E_{p_{\sigma(1)}},\dots,E_{p_{\sigma(k)}})\,
\varepsilon^J(E_{p_{\sigma(k+1)}},\dots,E_{p_{\sigma(k+l)}}).
\end{align*}
再由引理\ref{lemma:ch14-elementary-props}，上式和式中唯一给出非零值的项，是那些 $\sigma$ 分别单独排列 $P$ 的前 $k$ 个指标与后 $l$ 个指标的项。换言之，$\sigma$ 必形如 $\sigma=\tau\upsilon$，其中 $\tau\in S_k$ 作用于 $\{1,\dots,k\}$，$\upsilon\in S_l$ 作用于 $\{k+1,\dots,k+l\}$。由于 $\operatorname{sgn}(\tau\upsilon)=(\operatorname{sgn}\tau)(\operatorname{sgn}\upsilon)$，我们有
\begin{align*}
\varepsilon^I\wedge\varepsilon^J(E_{p_1},\dots,E_{p_{k+l}})
&=\frac1{k!\,l!}\sum_{\tau\in S_k}\sum_{\upsilon\in S_l}
(\operatorname{sgn}\tau)(\operatorname{sgn}\upsilon)\,
\varepsilon^I(E_{p_{\tau(1)}},\dots,E_{p_{\tau(k)}})\,
\varepsilon^J(E_{p_{k+\upsilon(1)}},\dots,E_{p_{k+\upsilon(l)}})\\
&=\left(\frac1{k!}\sum_{\tau\in S_k}(\operatorname{sgn}\tau)\,
\varepsilon^I(E_{p_{\tau(1)}},\dots,E_{p_{\tau(k)}})\right)
\left(\frac1{l!}\sum_{\upsilon\in S_l}(\operatorname{sgn}\upsilon)\,
\varepsilon^J(E_{p_{k+\upsilon(1)}},\dots,E_{p_{k+\upsilon(l)}})\right)\\
&=(\operatorname{Alt}\varepsilon^I)(E_{p_1},\dots,E_{p_k})\,
(\operatorname{Alt}\varepsilon^J)(E_{p_{k+1}},\dots,E_{p_{k+l}})\\
&=\varepsilon^I(E_{p_1},\dots,E_{p_k})\,
\varepsilon^J(E_{p_{k+1}},\dots,E_{p_{k+l}})=1.
\end{align*}

\emph{情形 4}：$P$ 是 $IJ$ 的一个排列且无重复指标。此时对 $P$ 作一个排列便回到情形~3。由于这个排列对\eqref{align:ch14-wedge-concat-eval}两边的作用是同乘一个相同的符号，结果在此情形下也成立。
\end{proof}

\begin{proposition}[外积的性质]\label{proposition:ch14-wedge-props}
设 $\omega,\omega',\eta,\eta',\xi$ 是有限维向量空间 $V$ 上的多重余向量。
\begin{enumerate}[(a)]
\item \emph{双线性}：对 $a,a'\in\mathbb{R}$，
\begin{align*}
(a\omega+a'\omega')\wedge\eta=a(\omega\wedge\eta)+a'(\omega'\wedge\eta),\\
\eta\wedge(a\omega+a'\omega')=a(\eta\wedge\omega)+a'(\eta\wedge\omega').
\end{align*}
\item \emph{结合性}：
\begin{align*}
\omega\wedge(\eta\wedge\xi)=(\omega\wedge\eta)\wedge\xi.
\end{align*}
\item \emph{反交换性}：对 $\omega\in\Lambda^k(V^*)$、$\eta\in\Lambda^l(V^*)$，
\begin{align}
\omega\wedge\eta=(-1)^{kl}\eta\wedge\omega.\label{align:ch14-anticomm}
\end{align}
\item 若 $(\varepsilon^i)$ 是 $V^*$ 的任意一组基，$I=(i_1,\dots,i_k)$ 是任意多重指标，则
\begin{align}
\varepsilon^{i_1}\wedge\cdots\wedge\varepsilon^{i_k}=\varepsilon^I.\label{align:ch14-epsilon-basis}
\end{align}
\item 对任意余向量 $\omega^1,\dots,\omega^k$ 与向量 $v_1,\dots,v_k$，
\begin{align}
\omega^1\wedge\cdots\wedge\omega^k(v_1,\dots,v_k)=\det\bigl(\omega^j(v_i)\bigr).\label{align:ch14-wedge-det}
\end{align}
\end{enumerate}
\end{proposition}

\begin{proof}
双线性直接来自定义，因为张量积是双线性的，而 $\operatorname{Alt}$ 是线性的。为证结合性，注意引理\ref{lemma:ch14-wedge-concat}给出
\begin{align*}
(\varepsilon^I\wedge\varepsilon^J)\wedge\varepsilon^K=\varepsilon^{IJ}\wedge\varepsilon^K
=\varepsilon^{IJK}=\varepsilon^I\wedge\varepsilon^{JK}
=\varepsilon^I\wedge(\varepsilon^J\wedge\varepsilon^K).
\end{align*}
一般情形由双线性得到。类似地，再次利用引理\ref{lemma:ch14-wedge-concat}，得
\begin{align*}
\varepsilon^I\wedge\varepsilon^J=\varepsilon^{IJ}=(\operatorname{sgn}\sigma)\varepsilon^{JI}
=(\operatorname{sgn}\sigma)\varepsilon^J\wedge\varepsilon^I,
\end{align*}
其中 $\sigma$ 是把 $IJ$ 变成 $JI$ 的排列。容易验证 $\operatorname{sgn}\sigma=(-1)^{kl}$，因为 $\sigma$ 可分解为 $kl$ 个对换的复合（$I$ 的每个指标都必须越过 $J$ 的每个指标）。反交换性随后由双线性得到。

(d) 是引理\ref{lemma:ch14-wedge-concat}与归纳的直接推论。为证 (e)，注意当每个 $\omega^j$ 都是基余向量 $\varepsilon^{i_j}$ 之一时，特殊情形归结为\eqref{align:ch14-epsilon-basis}。由于\eqref{align:ch14-wedge-det}两边都是 $(\omega^1,\dots,\omega^k)$ 的多重线性函数，这就足够了。
\end{proof}

由于该命题 (d) 部分的缘故，此后我们一般互换地使用记号 $\varepsilon^I$ 与 $\varepsilon^{i_1}\wedge\cdots\wedge\varepsilon^{i_k}$。

一个 $k$-余向量 $\eta$ 若能写成 $\eta=\omega^1\wedge\cdots\wedge\omega^k$（其中 $\omega^1,\dots,\omega^k$ 是余向量），就称为\emph{可分解的}。需要注意的是，当 $k>1$ 时并非每个 $k$-余向量都可分解（见习题\ref{exercise:ch14-02}）；然而由命题\ref{proposition:ch14-alternating-basis}与命题\ref{proposition:ch14-wedge-props}(d) 可知，每个 $k$-余向量都能写成可分解者的线性组合。

外积的定义与计算性质乍看可能令人望而生畏。然而，在大多数实际场合你只需记住上一命题所表达的那些性质。事实上，下面的练习表明，这些性质已足以唯一确定外积。

\begin{exercise}\label{exercise:ch14-12}
证明外积是唯一满足\eqref{align:ch14-epsilon-basis}的、结合的、双线性的、反交换的映射 $\Lambda^k(V^*)\times\Lambda^l(V^*)\to\Lambda^{k+l}(V^*)$。
\end{exercise}

对任意 $n$ 维向量空间 $V$，定义向量空间
\begin{align*}
\Lambda(V^*)=\bigoplus_{k=0}^n\Lambda^k(V^*).
\end{align*}
由命题\ref{proposition:ch14-alternating-basis}可知 $\dim\Lambda(V^*)=2^n$。命题\ref{proposition:ch14-wedge-props}表明，外积使 $\Lambda(V^*)$ 成为一个结合代数，称为 $V$ 的\emph{外代数}（或\emph{Grassmann 代数}）。这个代数不可交换，但它具有一个与之密切相关的性质。若一个代数 $A$ 有一个直和分解 $A=\bigoplus_{k\in\mathbb Z}A^k$，且乘积满足 $A^kA^l\subseteq A^{k+l}$，就称它是\emph{分次的}。若分次代数的乘积对 $a\in A^k$、$b\in A^l$ 满足 $ab=(-1)^{kl}ba$，就称它是\emph{反交换的}。命题\ref{proposition:ch14-wedge-props}(c) 表明 $\Lambda(V^*)$ 是一个反交换分次代数（其中我们约定 $A^k=\Lambda^k(V^*)$ 对 $0\leqslant k\leqslant n$，否则 $A^k=\{0\}$）。

正如我们在本节开头所观察到的，也可以不用\eqref{align:ch14-wedge-def}中那个笨重的系数来定义外积。有些作者采用这种替代定义。为避免混淆，我们记
\begin{align}
\omega\widetilde\wedge\eta=\operatorname{Alt}(\omega\otimes\eta).\label{align:ch14-alt-convention}
\end{align}
在这种定义下，\eqref{align:ch14-wedge-concat}换成
\begin{align*}
\varepsilon^I\widetilde\wedge\varepsilon^J=\frac{k!\,l!}{(k+l)!}\varepsilon^{IJ},
\end{align*}
而\eqref{align:ch14-wedge-det}换成
\begin{align}
\omega^1\widetilde\wedge\cdots\widetilde\wedge\omega^k(v_1,\dots,v_k)=\frac1{k!}\det\bigl(\omega^i(v_j)\bigr),\label{align:ch14-alt-det}
\end{align}
其中 $\omega^1,\dots,\omega^k$ 是余向量，你可以自行验证。

由于\eqref{align:ch14-wedge-det}，我们称\eqref{align:ch14-wedge-def}为外积的\emph{行列式约定}，而称\eqref{align:ch14-alt-convention}为\emph{Alt 约定}。采用哪种定义很大程度上是个人偏好问题。尽管 Alt 约定的定义或许略自然些，但行列式约定在计算上的优势使其在大多数应用中更可取，故本书一律采用它。（但习题\ref{exercise:ch14-03}给出了一个支持 Alt 约定的论证。）行列式约定在入门级微分几何教材中最常见，例如见于\cite{Boo86,Cha06,dC92,LeeJeff09,Pet06,Spi99}；Alt 约定见于\cite{KN69}，在复微分几何中更为常见。

\subsection{内乘}

有一个把向量与斜对称张量联系起来的重要运算。设 $V$ 是一个有限维向量空间。对每个 $v\in V$，我们定义一个线性映射 $i_v:\Lambda^k(V^*)\to\Lambda^{k-1}(V^*)$，称为由 $v$ 作的\emph{内乘}，定义如下：
\begin{align*}
i_v\omega(w_1,\dots,w_{k-1})=\omega(v,w_1,\dots,w_{k-1}).
\end{align*}
换言之，$i_v\omega$ 是把 $v$ 插入 $\omega$ 的第一个自变量位置而得到的。按惯例，当 $\omega$ 是一个 $0$-余向量（即一个数）时，我们把 $i_v\omega$ 理解为零。另一种常见记号是
\begin{align*}
v\lrcorner\omega=i_v\omega,
\end{align*}
常读作“$v$ 插入 $\omega$”。

\begin{lemma}\label{lemma:ch14-interior}
设 $V$ 是一个有限维向量空间，$v\in V$。
\begin{enumerate}[(a)]
\item $i_v\circ i_v=0$。
\item 若 $\omega\in\Lambda^k(V^*)$、$\eta\in\Lambda^l(V^*)$，则
\begin{align}
i_v(\omega\wedge\eta)=(i_v\omega)\wedge\eta+(-1)^k\omega\wedge(i_v\eta).\label{align:ch14-interior-deriv}
\end{align}
\end{enumerate}
\end{lemma}

\begin{proof}
对 $k\geqslant2$ 的 $k$-余向量，(a) 直接来自定义，因为任意斜对称张量在有两个自变量相同时取零值。对 $1$-余向量与 $0$-余向量，它由 $i_v$ 在 $0$-余向量上为零这一事实得到。

为证 (b)，只需考虑 $\omega$ 与 $\eta$ 都可分解的情形，因为任意正次数的斜对称张量都能写成可分解者的线性组合。在此特殊情形下，下面这个关于余向量 $\omega^1,\dots,\omega^k$ 的一般公式直接验证 (b)：
\begin{align}
v\lrcorner(\omega^1\wedge\cdots\wedge\omega^k)
=\sum_{i=1}^k(-1)^{i-1}\omega^i(v)\,\omega^1\wedge\cdots\wedge\widehat{\omega^i}\wedge\cdots\wedge\omega^k,\label{align:ch14-interior-expansion}
\end{align}
其中帽子表示删去 $\omega^i$ 这一项。

为证\eqref{align:ch14-interior-expansion}，令 $v_1=v$，把两边作用于任意 $(k-1)$ 向量组 $(v_2,\dots,v_k)$；于是我们要证的是
\begin{align}
(\omega^1\wedge\cdots\wedge\omega^k)(v_1,\dots,v_k)
=\sum_{i=1}^k(-1)^{i-1}\omega^i(v_1)\,
(\omega^1\wedge\cdots\wedge\widehat{\omega^i}\wedge\cdots\wedge\omega^k)(v_2,\dots,v_k).\label{align:ch14-interior-expand2}
\end{align}
\eqref{align:ch14-interior-expand2}的左端是矩阵 $\nu$ 的行列式，其 $(i,j)$ 元为 $\omega^i(v_j)$。为化简右端，令 $\nu^j_i$ 表示由 $\nu$ 删去第 $i$ 行、第 $j$ 列所得的 $(k-1)\times(k-1)$ 子矩阵。则\eqref{align:ch14-interior-expand2}的右端为
\begin{align*}
\sum_{i=1}^k(-1)^{i-1}\omega^i(v_1)\det\nu^1_i.
\end{align*}
这正是 $\det\nu$ 按第一列的余子式展开，因而等于 $\det\nu$。
\end{proof}

当外积采用 Alt 约定定义时，向量与一个 $k$-形式的内乘必须额外乘以一个因子 $k$ 来定义：
\begin{align*}
\widetilde{i}_v\omega(w_1,\dots,w_{k-1})=k\,\omega(v,w_1,\dots,w_{k-1}).
\end{align*}
这一定义确保内乘 $\widetilde{i}_v$ 仍满足\eqref{align:ch14-interior-deriv}——因子 $k$ 补偿了在 Alt 约定下计算\eqref{align:ch14-interior-expand2}左右两端时所出现的 $1/k!$ 与 $1/(k-1)!$ 之差。
\end{document}
